\section{Information Carriers}
\label{sec:information-carriers}

\subsection{The Spill Problem}

Suppose an intermediate value satisfies \(y\in[Y]\), where
\([Y]=\{0,1,\ldots,Y-1\}\). If \(Y\) is not a power of two, storing \(y\)
independently requires \(\lceil\log_2Y\rceil\) bits, although it contains only
\(\log_2Y\) bits of information. Repeating this upward rounding creates
linear redundancy.

The paper calls such a value a \emph{spill}. A spill is not lost information
and does not indicate a failed encoding; it is a reversible intermediate
state whose range has not yet been embedded into a standard binary memory
region.

\subsection{The Information Carrier Abstraction}

Let
\[
    x\in[X],
    \qquad
    y\in[Y].
\]
Here \(x\) is an information carrier and \(y\) is the old spill. The goal is
an injective local transformation
\[
    (x,y)\longmapsto(m,s),
\]
where
\[
    m\in[2^M],
    \qquad
    s\in[S].
\]
The value \(m\) is an \(M\)-bit memory content, while \(s\) is the new spill.

The asymmetry in the decoding requirement is essential: \(y\) must be
recoverable from \(m\) alone, whereas \(x\) may be recovered from \(m\) and
\(s\). Thus the old spill is truly committed to the current memory word. If
recovering it also required the new spill, queries could follow an unbounded
chain of successor spills and locality would be lost.

\subsection{The Information Carrier Lemma}

\begin{lemma}[Information carrier]
\label{lem:information-carrier}
Let \(x\in[X]\) and \(y\in[Y]\), with the parameters and values fitting in a
constant number of Word RAM words. There are integers \(M\) and \(S\), and an
injective mapping
\[
    [X]\times[Y]\longrightarrow[2^M]\times[S],
\]
such that \(y\) is recoverable from \(m\) alone, \(x\) is recoverable from
\((m,s)\), and
\[
    S=O(\sqrt X),
\]
\[
    2^M=O(Y\sqrt X),
\]
and
\[
    M+\log_2S-\log_2X-\log_2Y
    =O\left(\frac{1}{\sqrt X}\right).
\]
The mapping and its inverse take \(O(1)\) time.
\end{lemma}

\subsection{Construction}

Choose \(M\) so that \(2^M=\Theta(Y\sqrt X)\), with constants chosen to make
the following integer positive, and define
\[
    C=
    \left\lfloor
    \frac{2^M}{Y}
    \right\rfloor.
\]
Then \(C=\Theta(\sqrt X)\). Split the carrier into a quotient and remainder:
\[
    c=x\bmod C,
\]
\[
    s=\left\lfloor\frac{x}{C}\right\rfloor,
\]
and pack the old spill together with the remainder by setting
\[
    m=yC+c.
\]
The new spill has range size
\[
    S=
    \left\lceil
    \frac{X}{C}
    \right\rceil.
\]
Because \(c<C\) and \(y<Y\),
\[
    m<YC\leq2^M,
\]
so \(m\) is a valid \(M\)-bit memory content.

\subsection{Decoding and Correctness}

The old spill and the carrier remainder are obtained from \(m\) by
\[
    y=
    \left\lfloor
    \frac{m}{C}
    \right\rfloor,
\]
\[
    c=m\bmod C.
\]
In particular, \(m\) independently recovers \(y\). Given the new spill as
well, the carrier is reconstructed as
\[
    x=sC+c.
\]
These equations uniquely recover both inputs, proving that the mapping is
injective and that no information is discarded.

\subsection{Redundancy Analysis}

The redundancy has two independent rounding sources. First, the floor in
the definition of \(C\) leaves some of the \(2^M\) memory states unused when
\([Y]\times[C]\) is packed into memory. Its contribution is
\[
    R_1
    =M-\log_2(YC)
    =O\left(\frac{1}{C}\right).
\]
Indeed, \(2^M<Y(C+1)\), so the ratio between the available and used state
counts is less than \(1+1/C\), whose logarithm is \(O(1/C)\).

Second, the ceiling in
\(S=\lceil X/C\rceil\) may add an incomplete final quotient class when
\([X]\) is embedded into \([C]\times[S]\). This contribution is
\[
    R_2
    =\log_2 C
    +\log_2 S
    -\log_2 X
    =O\left(\frac{C}{X}\right).
\]
Here \(CS<X+C\), so the represented state count exceeds \(X\) by a factor
smaller than \(1+C/X\).

The total redundancy is therefore
\[
    R_1+R_2
    =O\left(
    \frac{1}{C}
    +
    \frac{C}{X}
    \right).
\]
Balancing the loss from floor rounding against the loss from ceiling rounding
gives
\[
    C=\Theta(\sqrt X),
\]
and hence
\[
    R_1+R_2
    =O\left(\frac{1}{\sqrt X}\right).
\]
The same choice yields \(S=O(\sqrt X)\) and
\(2^M=O(Y\sqrt X)\), completing the claimed bounds.

\subsection{Constant-Time Implementation}

Encoding and decoding use only a constant number of additions,
multiplications, divisions, and modulo operations. The range parameters are
fixed for an instance and can be precomputed. Under the paper's Word RAM
assumptions, these word-level operations take \(O(1)\) time. Division by a
fixed constant can also be replaced by multiplication with a precomputed
reciprocal and a constant amount of correction when appropriate.

\subsection{Why the Lemma Is Powerful}

The lemma is not an ordinary compression statement: it specifies which part
of the input becomes independently accessible. The old spill is committed
locally to \(m\), so decoding does not create an infinite dependence on later
carriers. Only the unresolved part of the current carrier continues as the
new spill. Moreover, increasing the carrier universe \(X\) reduces the
per-step redundancy as \(O(1/\sqrt X)\).

This mechanism is reused in the paper's vector representation and in its
online prefix-free code. The next application groups vector information into
large enough carriers that spills can be absorbed with negligible loss while
reads and updates remain local.
